\documentclass[11pt,a4paper]{article}

% Packages
\usepackage[utf8]{inputenc}
\usepackage[T1]{fontenc}
\usepackage[ngerman]{babel}
\usepackage{geometry}
\usepackage{graphicx}
\usepackage{xcolor}
\usepackage{booktabs}
\usepackage{fancyhdr}
\usepackage{tcolorbox}
\usepackage{enumitem}
\usepackage{hyperref}

% Page geometry
\geometry{margin=2.5cm}

% Colors (FehrAdvice style)
\definecolor{fehrblue}{RGB}{0,82,147}
\definecolor{fehrgray}{RGB}{100,100,100}
\definecolor{fehrlight}{RGB}{240,248,255}
\definecolor{fehrwarn}{RGB}{255,240,230}
\definecolor{fehrsuccess}{RGB}{230,255,230}

% Header/Footer
\pagestyle{fancy}
\fancyhf{}
\fancyhead[L]{\textcolor{fehrblue}{\textbf{EBF Management Report}}}
\fancyhead[R]{\textcolor{fehrgray}{Treppennutzung steigern}}
\fancyfoot[C]{\thepage}
\renewcommand{\headrulewidth}{0.5pt}

% Box styles
\tcbset{
    keymessage/.style={
        colback=fehrlight,
        colframe=fehrblue,
        fonttitle=\bfseries,
        title=#1,
        boxrule=0.5pt
    },
    warning/.style={
        colback=fehrwarn,
        colframe=orange,
        fonttitle=\bfseries,
        title=#1,
        boxrule=0.5pt
    },
    success/.style={
        colback=fehrsuccess,
        colframe=green!50!black,
        fonttitle=\bfseries,
        title=#1,
        boxrule=0.5pt
    }
}

\begin{document}

% Title Page
\begin{titlepage}
\centering
\vspace*{3cm}
{\Huge\textcolor{fehrblue}{\textbf{Treppe statt Lift}}\par}
\vspace{0.5cm}
{\Large\textcolor{fehrgray}{Management Report}\par}
\vspace{2cm}

\begin{tcolorbox}[keymessage={Verhaltensbasierte Intervention zur Gesundheitsförderung}]
Wie Mitarbeitende einer Schweizer Versicherung durch evidenzbasierte Nudges zur Treppennutzung motiviert werden können -- ohne Verbote, ohne Zwang.
\end{tcolorbox}

\vspace{2cm}

\begin{tabular}{ll}
\textbf{Kontext:} & Schweizer Versicherung \\
\textbf{Zielgruppe:} & Mitarbeitende (Bürogebäude) \\
\textbf{Modul:} & BCM2\_HLT\_STAIRS \\
\textbf{Version:} & 1.0 \\
\textbf{Datum:} & Januar 2026 \\
\end{tabular}

\vfill

{\large Evidence-Based Framework for Economic and Social Behavior\par}
\vspace{0.5cm}
{\textcolor{fehrgray}{FehrAdvice \& Partners AG}\par}
\end{titlepage}

% Executive Summary
\section*{\textcolor{fehrblue}{Executive Summary}}

\begin{tcolorbox}[keymessage={Kernbotschaft}]
\textbf{Menschen nehmen den Lift nicht aus Faulheit, sondern weil er der \textit{Default} ist.} Die Treppe ist unsichtbar, der Lift omnipräsent. Durch gezielte Choice Architecture -- Sichtbarkeit erhöhen, Defaults ändern, soziale Normen aktivieren -- kann die Treppennutzung um 30-50\% gesteigert werden.
\end{tcolorbox}

\subsection*{Das Problem}
\begin{itemize}[leftmargin=*]
    \item Nur 15-20\% der Mitarbeitenden nutzen regelmässig die Treppe
    \item Durchschnittlicher Büroangestellter: <2'000 Schritte/Tag im Gebäude
    \item Gesundheitskosten durch Bewegungsmangel: CHF 2.4 Mrd./Jahr (Schweiz)
    \item Versicherungen propagieren Gesundheit -- leben sie aber intern?
\end{itemize}

\subsection*{Die Lösung}
\begin{itemize}[leftmargin=*]
    \item \textbf{4 evidenzbasierte Interventionen} aus dem EBF-Toolkit
    \item \textbf{Kein Verbot} -- Lift bleibt verfügbar
    \item \textbf{Keine Kosten} für Infrastruktur (nur Kommunikation)
    \item \textbf{Messbare Wirkung} nach 4 Wochen
\end{itemize}

\subsection*{Erwarteter Impact}
\begin{itemize}[leftmargin=*]
    \item \textbf{+35\%} Treppennutzung (basierend auf Meta-Analyse, Nocon et al. 2010)
    \item \textbf{+1'200 Schritte/Tag} pro aktiven Mitarbeitenden
    \item \textbf{Signalwirkung:} Versicherung lebt, was sie predigt
\end{itemize}

\newpage

% Section 1: Verhaltensanalyse
\section{\textcolor{fehrblue}{Warum Menschen den Lift nehmen}}

Die Entscheidung Treppe vs. Lift ist keine bewusste Kosten-Nutzen-Abwägung. Sie geschieht in \textbf{System 1} -- automatisch, unbewusst, kontextgesteuert.

\subsection{Die vier Barrieren}

\begin{table}[h]
\centering
\begin{tabular}{llp{6cm}}
\toprule
\textbf{\#} & \textbf{Barriere} & \textbf{Erklärung} \\
\midrule
1 & \textbf{Unsichtbarkeit} & Treppe ist versteckt, Lift ist zentral \\
2 & \textbf{Default} & Lift ist der ``Normalfall'' \\
3 & \textbf{Soziale Norm} & Andere nehmen auch den Lift \\
4 & \textbf{Zeitperzeption} & Lift erscheint schneller (ist es meist nicht) \\
\bottomrule
\end{tabular}
\caption{Die vier verhaltensökonomischen Barrieren}
\end{table}

\begin{tcolorbox}[warning={Häufiger Fehler}]
\textbf{Information allein wirkt nicht.} Poster mit ``Treppe ist gesund!'' ändern nichts, wenn die Barrieren 1-4 bestehen bleiben. Die Intervention muss die \textit{Entscheidungsarchitektur} ändern, nicht das Wissen.
\end{tcolorbox}

\subsection{Der Schweizer Kontext ($\Psi_{CH}$)}

\begin{table}[h]
\centering
\begin{tabular}{lll}
\toprule
\textbf{Faktor} & \textbf{Wert} & \textbf{Implikation} \\
\midrule
Gesundheitsaffinität & Hoch (0.78) & Gesundheitsargumente resonieren \\
Eigenverantwortung & Hoch (0.82) & Keine Verbote, Wahlfreiheit betonen \\
Soziale Konformität & Mittel (0.55) & Normen wirken, aber subtil \\
Qualitätsbewusstsein & Sehr hoch (0.89) & ``Richtig machen'' als Argument \\
\bottomrule
\end{tabular}
\caption{Relevante Kontext-Parameter für die Schweiz (BCM2\_04\_KON)}
\end{table}

\newpage

% Section 2: Interventionen
\section{\textcolor{fehrblue}{Vier evidenzbasierte Interventionen}}

Alle Interventionen sind nach dem \textbf{20-Field Schema} (Appendix HHH) spezifiziert und auf den Schweizer Versicherungskontext kalibriert.

\subsection{Intervention 1: Point-of-Decision Prompts}

\begin{table}[h]
\centering
\begin{tabular}{ll}
\toprule
\textbf{Feld} & \textbf{Spezifikation} \\
\midrule
F1 Titel & ``Treppenschilder am Entscheidungspunkt'' \\
F2 Typ & Strukturell (Choice Architecture) \\
F3 Kontext & Organisation (Bürogebäude) \\
F5 FEPSDE & Physical (P) \\
F6 Journey-Phase & Action \\
\bottomrule
\end{tabular}
\end{table}

\textbf{Umsetzung:} Schilder direkt beim Lift mit Zeitvergleich:
\begin{quote}
\textit{``Treppe: 45 Sekunden bis Stock 3. Lift: 90 Sekunden (Wartezeit + Fahrt).''}
\end{quote}

\textbf{Evidenz:} Meta-Analyse zeigt +8-12\% Treppennutzung (Nocon et al. 2010)

\subsection{Intervention 2: Soziale Normen sichtbar machen}

\begin{table}[h]
\centering
\begin{tabular}{ll}
\toprule
\textbf{Feld} & \textbf{Spezifikation} \\
\midrule
F1 Titel & ``Treppennutzungs-Statistik'' \\
F2 Typ & Normativ (Social Proof) \\
F3 Kontext & Organisation \\
F5 FEPSDE & Social (S) \\
F6 Journey-Phase & Awareness $\rightarrow$ Trigger \\
\bottomrule
\end{tabular}
\end{table}

\textbf{Umsetzung:} Wöchentliche Anzeige im Eingangsbereich:
\begin{quote}
\textit{``Diese Woche: 847 Kolleginnen und Kollegen haben die Treppe genommen.''}
\end{quote}

\textbf{Evidenz:} Soziale Normen erhöhen Treppennutzung um +15-20\% (Burger \& Shelton, 2011)

\subsection{Intervention 3: Ästhetische Aufwertung}

\begin{table}[h]
\centering
\begin{tabular}{ll}
\toprule
\textbf{Feld} & \textbf{Spezifikation} \\
\midrule
F1 Titel & ``Treppen-Makeover'' \\
F2 Typ & Strukturell (Salienz) \\
F3 Kontext & Organisation \\
F5 FEPSDE & Emotional (E) + Physical (P) \\
F6 Journey-Phase & Maintenance \\
\bottomrule
\end{tabular}
\end{table}

\textbf{Umsetzung:}
\begin{itemize}
    \item Bessere Beleuchtung im Treppenhaus
    \item Motivierende Sprüche auf den Stufen (``Noch 12... 11... 10...'')
    \item Kunst oder Fotos an den Wänden
\end{itemize}

\textbf{Evidenz:} Ästhetische Aufwertung erhöht Nutzung um +10-15\% (Boutelle et al. 2001)

\subsection{Intervention 4: Gamification mit Team-Challenge}

\begin{table}[h]
\centering
\begin{tabular}{ll}
\toprule
\textbf{Feld} & \textbf{Spezifikation} \\
\midrule
F1 Titel & ``Stockwerk-Challenge'' \\
F2 Typ & Hybrid (Normativ + Feedback) \\
F3 Kontext & Organisation \\
F5 FEPSDE & Social (S) + Physical (P) \\
F6 Journey-Phase & Trigger $\rightarrow$ Maintenance \\
\bottomrule
\end{tabular}
\end{table}

\textbf{Umsetzung:} 4-wöchige Abteilungs-Challenge:
\begin{itemize}
    \item Teams sammeln ``Stockwerke'' (via App oder Strichliste)
    \item Wöchentliches Ranking
    \item Gewinnerteam: Symbolischer Preis (z.B. Team-Lunch)
\end{itemize}

\begin{tcolorbox}[warning={Crowding-Out Risiko}]
\textbf{Keine monetären Anreize!} Geldprämien würden die intrinsische Motivation (Gesundheit, Team) verdrängen ($\gamma_{Social,Financial} = -0.2$). Symbolische Anerkennung verstärkt die soziale Norm.
\end{tcolorbox}

\newpage

% Section 3: Implementierung
\section{\textcolor{fehrblue}{Implementierungsplan}}

\subsection{Portfolio-Logik}

Die vier Interventionen sind \textbf{komplementär} ($\gamma > 0$) und decken verschiedene Journey-Phasen ab:

\begin{table}[h]
\centering
\begin{tabular}{lcc}
\toprule
\textbf{Intervention} & \textbf{Phase} & \textbf{Komplementarität} \\
\midrule
1. Point-of-Decision & Action & -- \\
2. Soziale Normen & Awareness/Trigger & $\gamma_{1,2} = +0.3$ \\
3. Ästhetik & Maintenance & $\gamma_{1,3} = +0.2$ \\
4. Team-Challenge & Trigger/Maintenance & $\gamma_{2,4} = +0.4$ \\
\bottomrule
\end{tabular}
\caption{Komplementarität der Interventionen}
\end{table}

\subsection{Zeitplan}

\begin{table}[h]
\centering
\begin{tabular}{lll}
\toprule
\textbf{Woche} & \textbf{Aktion} & \textbf{Verantwortlich} \\
\midrule
0 & Baseline messen (Zähler installieren) & Facility Management \\
1 & Schilder aufhängen (Intervention 1) & Kommunikation \\
2 & Statistik-Display aktivieren (Int. 2) & HR/Kommunikation \\
3 & Treppen-Makeover (Intervention 3) & Facility Management \\
4-8 & Team-Challenge starten (Int. 4) & HR \\
9 & Evaluation \& Anpassung & Projektleitung \\
\bottomrule
\end{tabular}
\caption{Implementierungszeitplan}
\end{table}

\subsection{Erfolgsmessung}

\begin{itemize}[leftmargin=*]
    \item \textbf{Primär:} Treppennutzung (automatische Zähler)
    \item \textbf{Sekundär:} Mitarbeiterzufriedenheit (Pulse Survey)
    \item \textbf{Langfristig:} Krankheitstage, Gesundheitskosten
\end{itemize}

\begin{tcolorbox}[success={Erwartetes Ergebnis}]
\textbf{Prognose:} +35\% Treppennutzung nach 8 Wochen (95\% CI: [+25\%, +45\%])

Basierend auf: Meta-Analyse von 25 Studien (Nocon et al. 2010), adjustiert für Schweizer Kontext ($\Psi_{CH}$) und Versicherungsbranche.
\end{tcolorbox}

\newpage

% Quick Reference
\section*{\textcolor{fehrblue}{Quick Reference Card}}

\begin{tcolorbox}[keymessage={Checkliste für das Projektteam}]
\textbf{Vor dem Start:}
\begin{itemize}
    \item[$\square$] Baseline-Messung durchgeführt
    \item[$\square$] Facility Management informiert
    \item[$\square$] Budget für Schilder/Ästhetik genehmigt
    \item[$\square$] Challenge-App/Tool bereit
\end{itemize}

\textbf{Während der Umsetzung:}
\begin{itemize}
    \item[$\square$] Wöchentliche Datenerhebung
    \item[$\square$] Statistik-Display aktuell halten
    \item[$\square$] Challenge-Ranking kommunizieren
\end{itemize}

\textbf{Nach 8 Wochen:}
\begin{itemize}
    \item[$\square$] Evaluation durchführen
    \item[$\square$] Learnings dokumentieren
    \item[$\square$] Entscheidung: Fortführung/Anpassung
\end{itemize}
\end{tcolorbox}

\begin{tcolorbox}[keymessage={Die wichtigsten Prinzipien}]
\begin{enumerate}
    \item \textbf{Default ändern, nicht informieren} -- Der Lift ist das Problem, nicht das Wissen
    \item \textbf{Sichtbarkeit vor Motivation} -- Was man nicht sieht, wählt man nicht
    \item \textbf{Sozial vor monetär} -- Anerkennung wirkt stärker als Geld
    \item \textbf{Messen vor Meinen} -- Baseline $\rightarrow$ Intervention $\rightarrow$ Vergleich
\end{enumerate}
\end{tcolorbox}

\vfill

\begin{center}
\textit{Basierend auf EBF Appendix HHH (METHOD-TOOLKIT) und BCM2\_04\_KON (Schweizer Kontext)}

\vspace{0.5cm}

\textcopyright{} 2026 FehrAdvice \& Partners AG
\end{center}

\end{document}
